\documentclass[12pt]{article}

\setlength{\parindent}{0pt}
\setcounter{secnumdepth}{3}
\usepackage[margin=1in]{geometry}

% packages
\usepackage{amsmath, amsfonts, amsthm, amssymb}
\usepackage{graphicx, float}
\usepackage{mathtools}
\usepackage{titlesec}
\usepackage{interval}
\usepackage{hyperref}
\usepackage{siunitx}
\usepackage{titling}
\usepackage{vwcol}
\usepackage{setspace}
\usepackage{empheq}
\usepackage{cancel}
\usepackage{esdiff}
\usepackage{multicol}
\usepackage{mdframed}
\usepackage{esdiff}
\usepackage{tikzsymbols}
\usepackage{multicol}
\usepackage{tikz}
\usepackage{varwidth}
\usepackage{pgfplots}
\usepackage{fancyhdr}
\usepackage{enumitem}
\usepackage{tcolorbox}

\intervalconfig {
	soft open fences
}

% header
\pagestyle{fancy}
\fancyhf{}
\lhead{Ethan Fan}
\rhead{AP Calculus BC}
\cfoot{\thepage}

% section format
\titleformat{\section}
  {\large \bfseries}
  {}
  {0em}
  {}
\titlespacing*{\section}{0pt}{0.5em}{0.3em}

\titleformat{\subsection}[runin]
  {\bfseries}
  {}
  {0em}
  {}
\titlespacing*{\subsection}{0pt}{0pt}{0em}

\titleformat{\subsubsection}[runin]
  {\itshape}
  {(\roman{subsubsection})}
  {0em}
  {}

% commands
\newcommand{\alignedintertext}[1]{%
  \noalign{%
    \vskip\belowdisplayshortskip
    \vtop{\hsize=\linewidth#1\par
    \expandafter}%
    \expandafter\prevdepth\the\prevdepth
  }%
}

\newcommand*{\paren}[1]{\ensuremath\left(#1\right)}
\newcommand*{\limit}[2][x]{\ensuremath{\displaystyle\lim_{#1\to#2}}}
\newcommand*{\Limit}[3][x]{\ensuremath{\displaystyle\lim_{#1\to#2}\left[#3\right]}}
\newcommand*{\deriv}[1][x]{\ensuremath{\dfrac{\mathrm{d}}{\mathrm{d}#1}}}
\newcommand*{\Deriv}[2][x]{\ensuremath{\dfrac{\mathrm{d}}{\mathrm{d}#1}\left[#2\right]}}
\newcommand{\ddx}{\dfrac{d}{dx}}
\newcommand{\dydx}{\dfrac{dy}{dx}}
\newcommand{\dydt}{\dfrac{dy}{dt}}
\newcommand{\dxdt}{\dfrac{dx}{dt}}
\newcommand*{\iinteg}[2][x]{\ensuremath{\displaystyle\int #2\;\mathrm{d}#1}}
\newcommand*{\dinteg}[4][x]{\ensuremath{\displaystyle\int_{#2}^{#3}#4\;\mathrm{d}#1}}
\newcommand*{\abs}[1]{\ensuremath{\left|#1\right|}}
\newcommand*{\eps}{\varepsilon}
\newcommand*{\real}{\ensuremath\mathbb{R}}
\newcommand*{\integer}{\ensuremath\mathbb{Z}}
\newcommand*{\posint}{\ensuremath\mathbb{Z^+}}
\newcommand*{\cbrt}[1]{\ensuremath{\sqrt[3]{#1}}}
\newcommand*{\lheqzero}{\ensuremath{\underset{\text{L'H}}{\overset{\left[\frac00\right]}{=}}}}
\newcommand*{\lheqinfty}{\ensuremath{\underset{\text{L'H}}{\overset{\left[\frac{\infty}{\infty}\right]}{=}}}}
\newcommand{\tor}{\text{ or }}
\newcommand{\floor}[1]{\lfloor #1 \rfloor}
\newcommand{\at}{\Big|}

\newcommand{\vem}{\vspace{0.5em}}
\newcommand{\example}[1]{\section{Example #1}}
\newcommand{\problem}[1]{\section{Problem #1}}
\newcommand{\subproblem}[1]{\subsection{(#1)} \,}

\DeclareMathOperator{\dne}{DNE}
\DeclareMathOperator{\sgn}{sgn}

\newcommand{\definition}[1]{\begin{tcolorbox}[colback=red!5!white,colframe=red!75!black,parbox=false] #1 \end{tcolorbox}}

\newcommand{\theorem}[2]{\begin{tcolorbox}[title={#1},colback=blue!5!white,colframe=blue!75!black,parbox=false] #2 \end{tcolorbox}}

\allowdisplaybreaks
% \postdisplaypenalty=100000

% document
\begin{document}

\vspace*{0cm}

% opening
\begin{center}
    {\Large \bfseries Problem Set 25} \\[0.5cm]
    {\large Area of a Surface of Revolution} \\[0.35cm]
    {\normalsize November 9, 2026}
\end{center}

% problems

\problem{1}

\subproblem{a}
\begin{gather*}
    \frac{r_1}{l_1}=\frac{r_2}{l_1+l}\\
    r_1 l_1 + r_1 l = r_2 l_1\\
    \boxed{(r_2-r_1)l_1=r_1l}\\
    A=\pi r_2(l_1+l)-\pi r_1l_1=\boxed{\pi (r_1+r_2)l}
\end{gather*}

\subproblem{b}
\begin{gather*}
    l=\sqrt{15^2+(18-10)^2}=17\\
    A=\pi (10+18)17=\boxed{476\pi \text{ cm}^2}
\end{gather*}

\subproblem{c}
\begin{gather*}
    8\cdot l_1=10\cdot 17\implies l_1=85/4\\
    85/4 \cdot \theta =2\pi \cdot 10\\
    R_2=17+85/4=153/4\\
    \theta =\boxed{\frac{16\pi}{17}}
\end{gather*}
The angle is less than $180^\circ$. 

\problem{2}

\subproblem{c}
\begin{gather*}
    y=e^x \implies y'=e^x\\
    A=\int_0^12\pi e^x\sqrt{1+e^{2x}} \, dx\\
    = \begin{bmatrix}
        u=e^{x}\\
        du=e^x\, dx
    \end{bmatrix} 2\pi \int_1^e \sqrt{1+u^2} \, du\\
    =2\pi (\frac{u}{2}\sqrt{1+u^2}+\frac{1}{2}\ln (u+\sqrt{1+u^2}) \at_1^e\\
    =\boxed{ \pi \paren{ e\sqrt{1+e^2} - \sqrt{ 2}+ \ln \paren{ \frac{e+\sqrt{1+e^2}}{1+\sqrt{2}}}}}
\end{gather*}

\problem{3}

\subproblem{b}
\begin{gather*}
    y=\sin \pi x\implies y' = \pi \cos \pi x\\
    A=2\pi\int _0^1 \sin \pi x \sqrt{1+\pi^2\cos^2 \pi x} \, dx\\
    = \begin{bmatrix}
        u=\pi \cos \pi x\\
        du=- \pi^2 \sin x\, dx
    \end{bmatrix}
    \frac{2}{\pi} \int_{-\pi}^\pi \sqrt{1+u^2 } \, du\\
    =\frac{4}{\pi} (\frac{u}{2}\sqrt{1+u^2}+\frac{1}{2}\ln (u+\sqrt{1+u^2}) \at_0^\pi \\
    = \boxed{2\sqrt{1+\pi^2} + \frac{2}{\pi} \ln (\pi + \sqrt{1+\pi^2})}
\end{gather*}

\subproblem{d}
\begin{gather*}
    x=y^3 \implies x'=3y^2\\
    A=2\pi \int_1^2 y^3\sqrt{1+9y^4} \, dy\\
    = \begin{bmatrix}
        u=y^4\\
        du=3y^3\, dy
    \end{bmatrix} \frac{2}{3}\int_1^{16} \sqrt{1+9y} \, dy\\
    =  \frac{2}{3} \cdot \frac{1}{9} \cdot \frac{2}{3} (1+9y)^\frac{3}{2}\at_1^{16}\\
    = \boxed{\frac{5\pi}{27} \paren{29\sqrt{145} - 2\sqrt{10}}}
\end{gather*}

\problem{4}

\subproblem{d}
\begin{gather*}
    H=R-R\sin \varphi\\
    A=\boxed{\frac{1}{2}(1-\sin \varphi)}\\
    2\cdot \sin 30^\circ \cdot \frac{1}{2} = \boxed{\frac{1}{2}}
\end{gather*}

\subproblem{e}
\begin{gather*}
    \cos \alpha = R/R+h \implies \sin \frac{\pi}{2}-\alpha =R/R+h\\
    A=\frac{1}{2}(\frac{R+h-R}{R+h})= \boxed{\frac{h}{2(R+h)}}
\end{gather*}
No satellite can see half of the Earth. 

\problem{5}

\subproblem{a}
\begin{gather*}
    y=\frac{1}{2}x^2\implies y'=x\\
    A=2\pi \int_0^\frac{1}{2} x\sqrt{1+x^2} \, dx\\
    = \begin{bmatrix}
        u=x^2\\
        du=2x\,d x
    \end{bmatrix} \int_0^\frac{1}{4} \sqrt{1+u} \, du\\
    = \frac{2}{3}u^\frac{3}{2} \at_1^\frac{5}{4}\\
    =\frac{\pi(5\sqrt{5} - 8)}{12} \text{ m}^2\\
\end{gather*}

\subproblem{b}
\begin{gather*}
    ds=\sqrt{1+(y')^2} \, dx >dx
\end{gather*}

\problem{6}

\subproblem{a}
\begin{gather*}
    A=2\pi \int_C (r-y)\, ds\\
    = 2\pi \int_C r\, ds -2\pi \int_C y\, ds\\
    =\boxed{4\pi^2r^2}
\end{gather*}

\subproblem{b}
\begin{gather*}
    \overline{y} = \frac{1}{L} \int_C y\, ds\\
    S=2\pi \int_Cy\, ds=\boxed{2\pi \overline{y}L}
\end{gather*}

\subproblem{c}
\begin{gather*}
    S=2\pi \cdot \frac{r}{2}\cdot L =\boxed{\pi rl}\\
    S=2\pi \cdot 3\pi \cdot 4=\boxed{24\pi^2 \text{ cm}^2}
\end{gather*}

\problem{7}

\subproblem{a}
\begin{gather*}
    A=\pi \int_1^b \frac{1}{x^2} \, dx\\
    = \pi (-\frac{1}{x})\at_1^b\\
    = \pi - \frac{\pi}{b} <\pi
\end{gather*}

\subproblem{b}
\begin{gather*}
    y=\frac{1}{x} \implies y'=- \frac{1}{x^2}\\
    S(b)=2\pi \int_1^b \frac{1}{x} \sqrt{1+\frac{1}{x^4}} \, dx\\
    \frac{1}{x}\sqrt{1+\frac{1}{x^4}} > \frac{1}{x}\\
    S(b) >2\pi \int_1^b \frac{1}{x}  \, dx = 2\pi \ln b\\
    \lim_{b \to \infty} S(b) \to \infty
\end{gather*}

\subproblem{c}
The surface area of the horn can extend to infinity, while the volume will ultimately converge as the thickness of the horn approaches zero. 

\problem{8}

\subproblem{a}
\begin{gather*}
    y=\frac{x^2}{4}-\frac{\ln x}{2} \implies y'= \frac{x}{2}- \frac{1}{2x}\\
    \sqrt{1+(y')^2} = \sqrt{(\frac{x}{2}+\frac{1}{2x})^2} = \boxed{\frac{x}{2}+\frac{1}{2x}}\\
    l= \int_1^e \frac{x}{2}+\frac{1}{2x} = \frac{x^2}{4}+\frac{1}{2} \ln x \at_1^e= \boxed{\frac{e^2+1}{4}}
\end{gather*}

\subproblem{b}
\begin{gather*}
    S=2\pi \int_1^e x(\frac{x}{2}+\frac{1}{2x})\, dx\\
    = \pi \int_1^e x^2+1\, dx\\
    = \pi (\frac{x^3}{3}+x)\at_1^e\\
    = \boxed{\frac{\pi(e^3+3e-4)}{3}}
\end{gather*}

\subproblem{c}
\begin{gather*}
    S=2\pi \int_1^e (\frac{x^2}{4}-\frac{\ln x}{2})(\frac{x}{2}+\frac{1}{2x})\, dx\\
    = \frac{\pi}{2} \int_1^e \frac{x^3}{2}+\frac{x}{2}-x\ln x-\frac{\ln x}{x}\, dx\\
    =\boxed{\frac{\pi(e^4-9)}{16}}
\end{gather*}

\subproblem{d}
\begin{gather*}
    V = 2\pi \int_1^e x(\frac{x^2}{4}-\frac{\ln x}{2}) \, dx\\
    = \pi  \int_1^e \frac{x^3}{2}- x\ln x \, dx\\
    =\boxed{\frac{\pi (e^4-2e^2-3)}{8}}
\end{gather*}

\subproblem{e}
\[
ds>dx
\]

\problem{9}

\subproblem{a}
\begin{gather*}
    c=\frac{e^x+e^{-x}}{2} \implies c'=\frac{e^x-e^{-x}}{2}\\
    1+(c')^2 x = 1+\frac{e^{2x}}{4} -\frac{1}{2} + \frac{e^{-2x}}{4}= c^2
\end{gather*}

\subproblem{b}
\begin{gather*}
    S= 2\pi \int_a^b \paren{\frac{e^x+e^{-x}}{2}}^2 \, dx\\
    V= \pi \int_a^b\paren{\frac{e^x+e^{-x}}{2}}^2 \, dx\\
    \boxed{S=2V}
\end{gather*}

\subproblem{c}
\begin{gather*}
    S=2v\\
    2\pi \int_{a}^{b} f(x)\sqrt{1 + f'(x)^2} \, dx = 2\pi \int_{a}^{b} f(x)^2 \, dx\\
    \sqrt{1 + f'(x)^2}=f(x)\\
    \boxed{1 + f'(x)^2=f(x)^2}\\
    2f'(x) f''(x)=2f(x)f'(x)\\
    f'(x) \ne 0 \implies \boxed{f''(x)=f(x)}
\end{gather*}

\subproblem{d}
\begin{gather*}
    S= 2\pi \int_{-1}^1 \paren{\frac{e^x+e^{-x}}{2}}^2 \, dx \\
    = \frac{\pi}{2} \int_{-1}^1 e^{2x}+2+e^{-2x} \, dx\\
    =\frac{\pi}{2}(\frac{e^{2x}}{2}+2x-\frac{e^{-2x}}{2})\at_{-1}^1\\
    = \boxed{\pi \paren { \frac{e^2-e^{-2}}{2}+2} \approx 17.68}\\
    A=2\pi \cdot 2\cdot \frac{e^x+e^{-x}}{2}\\
    = \boxed{2\pi (e+\frac{1}{e}) \approx 19.39}\\
    A=\boxed{2\pi \paren{\frac{e+e^{-1}}{2}}^2 \approx 14.96}\\
    \boxed{14.96<17.68<19.39}
\end{gather*}   



































\end{document}

